\documentclass{article}
\usepackage[utf8]{inputenc}
\usepackage{amsmath}
\usepackage{amssymb}
\usepackage{enumitem}
\usepackage{graphicx}

\begin{document}

% Portada
\begin{titlepage}
    \centering
    
\includegraphics[height=2cm]{fing.png}\hfill
    
\includegraphics[height=2cm]{uach.png}
    \\
    \vspace{1.5cm}
    {\Large \textbf{Universidad Autónoma de Chihuahua}}\\[0.2cm]
    {\large Facultad de Ingeniería}\\[0.5cm]
    {\Large \textbf{Métodos Numéricos}}\\[0.5cm]
    {\large Examen Parcialito 1}\\[1cm]

    \vspace{1cm}
    {\large Nombre: \underline{Francisco Javier Ponce Saenz}}\\[0.3cm]
    {\large Matrícula: \underline{a325000}}\\[0.3cm]
    {\large Semestre: \underline{3er Semestre}}\\[0.3cm]
    {\large Fecha: \underline{\hspace{4cm}}}\\
    \vfill
\end{titlepage}

\section*{Examen}

Obtener la raíz de las siguientes funciones por el método que se indique:

\begin{enumerate}[label=\alph*)]
    \item Método Newton--Raphson mejorado:
    \[
    \text{Tol}: 0.001,\quad y=\sin(x),\quad x_0=0.6
    \]
    Respuesta:\rule{3cm}{0.4pt}
    
    \begin{figure}[h]
        \centering
        % tabla adicional para el método de Newton--Raphson
        \includegraphics[width=0.6\textwidth]{tabla1.png}
        \caption{Tabla adicional para el método de Newton--Raphson.}
    \end{figure}
    
    \item Método de falsa posición:
    \[
    y = \frac{\sqrt{x^2+2}\,(x+5)}{(x^2-3x-2)^2\,(x+4)^3},\quad a=0,\ b=7,\ \text{Tol}=0.01
    \]
    Respuesta:\rule{3cm}{0.4pt}
    
    \begin{figure}[h]
        \centering
        % tabla para falsa posición
        \includegraphics[width=0.6\textwidth]{tabla2.png}
        \caption{Tabla adicional al aplicar falsa posición.}
    \end{figure}
    \noindent Debido a que $f(0)=113.137$ y $f(7)=168.732$ tienen el mismo signo, el teorema del valor intermedio no
    garantiza ninguna raíz en $[0,7]$, por lo que el método falla.
    
    \item Método de bisección del intervalo:
    \[
    y = \frac{\left(\sqrt{3x^3+2x^2+4}\right)^2\,(x-3)}{(x^2+6x+7)^1\,(x^2+7)^2},\quad a=0,\ b=5,\ \text{Tol}=0.01
    \]
    Respuesta:\rule{3cm}{0.4pt}
    \\
    \begin{figure}[h]
        \centering
        % tabla para bisección
        \includegraphics[width=0.6\textwidth]{tabla3.png}
        \caption{Tabla adicional para el método de bisección.}
    \end{figure}
\end{enumerate}

\end{document}
